\documentclass[11pt,a4paper]{article}

\usepackage[utf8]{inputenc}
\usepackage[margin=0.9in]{geometry}
\usepackage{amsmath,amssymb,amsthm,amsfonts}
\usepackage{mathtools}
\usepackage{hyperref}
\usepackage{booktabs}
\usepackage{cite}
\usepackage{microtype}
\usepackage{array}

\hypersetup{
    colorlinks=true,
    linkcolor=blue,
    citecolor=red,
    urlcolor=blue
}

% Theorem Environments
\theoremstyle{plain}
\newtheorem{theorem}{Theorem}[section]
\newtheorem{lemma}[theorem]{Lemma}
\newtheorem{proposition}[theorem]{Proposition}
\newtheorem{corollary}[theorem]{Corollary}

\theoremstyle{definition}
\newtheorem{definition}[theorem]{Definition}
\newtheorem{example}[theorem]{Example}

\theoremstyle{remark}
\newtheorem{remark}[theorem]{Remark}

\numberwithin{equation}{section}

\title{\textbf{Nonlinear Simplicial Waves and Anomalous Porous Transport Induced by Beta-Kernel Fractional Laplacians}}
\author{\textbf{Reinaldo Maia Silva-Filho} \\ \textit{Universidade Federal de Lavras (UFLA)}}
\date{\today}


\DeclareMathOperator*{\argmin}{argmin}
\providecommand{\doi}[1]{\href{https://doi.org/#1}{\texttt{doi:#1}}}
\providecommand{\Hyp}{\mathbb{H}}
\providecommand{\Sph}{\mathbb{S}}
\providecommand{\II}{\mathrm{II}}
\providecommand{\R}{\mathbb{R}}
\providecommand{\Area}{\mathrm{Area}}
\providecommand{\Hsn}{\mathcal{H}}
\providecommand{\BV}{\mathrm{BV}}
\providecommand{\TV}{\mathrm{TV}}
\providecommand{\cutnorm}[1]{\|#1\|_{\square}}

\begin{document}

\maketitle

\begin{abstract}
We develop a comprehensive analytical and mathematical physics framework for non-local wave dynamics and subdiffusive transport governed by the \emph{Fractional Simplicial Laplacian} $\Delta_{\Delta_m}^\alpha$, an operator induced by continuous multinomial Beta-kernels on the standard simplex $\Delta_{m-1}(\alpha)$. In the first part of this work, we formulate the \emph{Nonlinear Simplicial Schr\"odinger Equation} (NLSE). We prove the global conservation of particle mass $\mathcal{N}[\psi]$ and simplicial Hamiltonian energy $\mathcal{E}[\psi]$ via the self-adjointness of the continuous Beta-integral kernel. We perform a modulational instability analysis demonstrating that plane waves undergo directional sideband destabilization governed by the discrete $S_m$ point-group symmetry of the simplex, producing facet-aligned anisotropic solitary waves. In the second part, we formulate the space-time fractional simplicial diffusion equation describing anomalous transport in anisotropic porous, fractured, and crystalline media. Using Fourier--Laplace transform methods, we derive the exact non-local propagator in terms of Mittag-Leffler functions and rigorously prove that the long-time Mean Squared Displacement (MSD) covariance tensor scales as $\langle \mathbf{x}\mathbf{x}^T \rangle(t) = \frac{\mathcal{K}_{\text{diff}}}{m \alpha \Gamma(\beta+1)} \mathbf{A}_{m-1} t^\beta$, where $\mathbf{A}_{m-1}$ is the Cartan matrix of the Lie algebra $A_{m-1}$. This establishes a rigorous geometric connection between Lie algebraic root systems, compact-support memory kernels, and anisotropic anomalous diffusion in structured media.
\end{abstract}

\tableofcontents
\newpage

% ==============================================================================
\section{Introduction and Physical Motivation}
% ==============================================================================

Non-local partial differential equations (PDEs) and fractional calculus have become foundational tools for describing memory-dependent dynamics, anomalous transport, and non-local wave phenomena across quantum physics, fluid mechanics, and geophysical transport \cite{podlubny1998, samko1993, metzler2000, laskin2000}. In traditional fractional quantum mechanics and anomalous diffusion models, the non-local spatial operator is typically chosen as the isotropic Riesz fractional Laplacian $(-\Delta)^{\alpha/2}$, whose integral representation:
\begin{equation}
(-\Delta)^{\alpha/2} u(\mathbf{x}) = C_{d, \alpha} \operatorname{P.V.} \int_{\mathbb{R}^d} \frac{u(\mathbf{x}) - u(\mathbf{y})}{|\mathbf{x} - \mathbf{y}|^{d+\alpha}} \, d\mathbf{y},
\end{equation}
is isotropic, radially symmetric under the orthogonal group $O(d)$, and characterized by algebraic power-law tails that extend across the entire space $\mathbb{R}^d$.

However, in many realistic physical media—such as anisotropic fractured rocks, crystalline waveguides, packed colloidal arrays, and simplicial metamaterials—the underlying microstructure imposes two fundamental physical constraints:
\begin{enumerate}
    \item \textbf{Directional Geometric Anisotropy:} Transport occurs predominantly along structured interconnected channels with characteristic intersection angles (such as triangular or tetrahedral lattices), breaking continuous rotational symmetry.
    \item \textbf{Finite Horizon of Non-Local Steps:} Microscopic jump lengths per collision event are geometrically bounded by the pore or cell size, precluding unphysical infinite-velocity jump tails.
\end{enumerate}

To bridge this fundamental gap, we introduce and analyze dynamical and transport processes governed by the \emph{Fractional Simplicial Laplacian} $\Delta_{\Delta_m}^\alpha$, constructed from the continuous multinomial Beta-kernel on the standard $(m-1)$-simplex $\Delta_{m-1}(\alpha)$.

% ==============================================================================
\section{The Fractional Simplicial Laplacian and its Metric Geometry}
% ==============================================================================

Let $\mathbf{x} = (x_1, \dots, x_{m-1}) \in \mathbb{R}^{m-1}$ denote spatial coordinates, and let $\Delta_{m-1}(\alpha) \subset \mathbb{R}_{\ge 0}^{m-1}$ be the standard scaled simplex of size $\alpha > 0$:
\begin{equation}
\Delta_{m-1}(\alpha) \coloneqq \left\{\mathbf{y} \in \mathbb{R}_{\ge 0}^{m-1} \;\middle|\; \sum_{j=1}^{m-1} y_j \le \alpha\right\}.
\end{equation}
The dependent barycentric coordinate is defined by $y_m \coloneqq \alpha - \sum_{j=1}^{m-1} y_j$. The continuous multinomial Beta-kernel is defined by
\begin{equation}
\binom{\alpha}{\mathbf{y}} \coloneqq \frac{\Gamma(\alpha+1)}{\prod_{j=1}^m \Gamma(y_j+1)}, \quad I_m(\alpha) \coloneqq \int_{\Delta_{m-1}(\alpha)} \binom{\alpha}{\mathbf{y}} \, d\mathbf{y}.
\end{equation}

\begin{definition}[Fractional Simplicial Laplacian]
For fractional parameter $\alpha > 0$, the self-adjoint Fractional Simplicial Laplacian acting on $u \in H^2(\mathbb{R}^{m-1})$ is defined by
\begin{equation}
\label{eq:simplicial_laplacian_def}
\Delta_{\Delta_m}^\alpha u(\mathbf{x}) \coloneqq \frac{1}{2\alpha^2 I_m(\alpha)}\int_{\Delta_{m-1}(\alpha)} \binom{\alpha}{\mathbf{y}} \Big[u(\mathbf{x}+\mathbf{y}) - 2u(\mathbf{x}) + u(\mathbf{x}-\mathbf{y})\Big] d\mathbf{y}.
\end{equation}
\end{definition}

\begin{proposition}[Self-Adjointness and Positivity]
The operator $-\Delta_{\Delta_m}^\alpha$ is linear, symmetric, and positive semi-definite on $L^2(\mathbb{R}^{m-1})$:
\begin{equation}
\langle u, -\Delta_{\Delta_m}^\alpha u \rangle_{L^2} = \frac{1}{2\alpha^2 I_m(\alpha)}\int_{\mathbb{R}^{m-1}} d\mathbf{x} \int_{\Delta_{m-1}(\alpha)} \binom{\alpha}{\mathbf{y}} |u(\mathbf{x}+\mathbf{y}) - u(\mathbf{x})|^2 d\mathbf{y} \ge 0.
\end{equation}
\end{proposition}

\begin{proof}
Expanding the inner product $\langle u, -\Delta_{\Delta_m}^\alpha u \rangle_{L^2} = \int_{\mathbb{R}^{m-1}} u^*(\mathbf{x}) (-\Delta_{\Delta_m}^\alpha u(\mathbf{x})) d\mathbf{x}$:
\begin{align*}
\langle u, -\Delta_{\Delta_m}^\alpha u \rangle &= \frac{1}{2\alpha^2 I_m(\alpha)}\int_{\Delta_{m-1}(\alpha)} \binom{\alpha}{\mathbf{y}} \notag \\
&\quad \times \left[ \int_{\mathbb{R}^{m-1}} \Big(2|u(\mathbf{x})|^2 - u^*(\mathbf{x})u(\mathbf{x}+\mathbf{y}) - u^*(\mathbf{x})u(\mathbf{x}-\mathbf{y})\Big) d\mathbf{x} \right] d\mathbf{y}.
\end{align*}
Using the translation invariance $\int u^*(\mathbf{x})u(\mathbf{x}-\mathbf{y})d\mathbf{x} = \int u^*(\mathbf{x}+\mathbf{y})u(\mathbf{x})d\mathbf{x}$:
\begin{equation}
\int_{\mathbb{R}^{m-1}} \Big(2|u(\mathbf{x})|^2 - u^*(\mathbf{x})u(\mathbf{x}+\mathbf{y}) - u(\mathbf{x})u^*(\mathbf{x}+\mathbf{y})\Big) d\mathbf{x} = \int_{\mathbb{R}^{m-1}} |u(\mathbf{x}+\mathbf{y}) - u(\mathbf{x})|^2 d\mathbf{x}.
\end{equation}
Since $\binom{\alpha}{\mathbf{y}} > 0$ on the interior of $\Delta_{m-1}(\alpha)$, the integral is strictly non-negative, vanishing if and only if $u$ is constant.
\end{proof}

\begin{theorem}[Dispersion Symbol and Cartan Metric Emergence]
\label{thm:dispersion}
Let $\mathbf{k} = (k_1, \dots, k_{m-1}) \in \mathbb{R}^{m-1}$ with the affine simplicial projection $k_m \equiv -\sum_{j=1}^{m-1} k_j$ (enforcing $\sum_{j=1}^m k_j = 0$ on the simplex root hyperplane). The Fourier symbol $\sigma_{\Delta_m}^\alpha(\mathbf{k})$ defined by $-\Delta_{\Delta_m}^\alpha e^{i\mathbf{k}\cdot\mathbf{x}} = \sigma_{\Delta_m}^\alpha(\mathbf{k}) e^{i\mathbf{k}\cdot\mathbf{x}}$ is given explicitly by
\begin{equation}
\sigma_{\Delta_m}^\alpha(\mathbf{k}) = \frac{1}{\alpha^2}\left[1 - R(\mathbf{k})^\alpha \cos\big(\alpha \Theta(\mathbf{k})\big)\right],
\end{equation}
where $R(\mathbf{k}) \coloneqq \frac{1}{m}\left|\sum_{j=1}^m e^{-ik_j}\right| \le 1$ and $\Theta(\mathbf{k}) \coloneqq \operatorname{Arg}\left(\sum_{j=1}^m e^{-ik_j}\right)$.

In the long-wavelength limit $|\mathbf{k}| \to 0$:
\begin{equation}
\lim_{|\mathbf{k}| \to 0} \sigma_{\Delta_m}^\alpha(\mathbf{k}) = \frac{1}{2m\alpha}\mathbf{k}^T \mathbf{A}_{m-1}\mathbf{k},
\end{equation}
where $\mathbf{A}_{m-1}$ is the $(m-1)\times(m-1)$ Cartan matrix of the simple Lie algebra $A_{m-1}$:
\begin{equation}
\mathbf{A}_{m-1} = \begin{pmatrix}
2 & 1 & 1 & \cdots & 1 \\
1 & 2 & 1 & \cdots & 1 \\
\vdots & \vdots & \vdots & \ddots & \vdots \\
1 & 1 & 1 & \cdots & 2
\end{pmatrix}.
\end{equation}
\end{theorem}

% ==============================================================================
\section{Nonlinear Simplicial Schr\"odinger Equations (NLSE)}
% ==============================================================================

We formulate the Cauchy problem for the non-local Fractional Simplicial NLSE on $\mathbb{R}^{m-1} \times \mathbb{R}^+$:
\begin{equation}
\label{eq:nlse_main}
i \hbar \frac{\partial \psi(\mathbf{x}, t)}{\partial t} = -\frac{\hbar^2}{2M}\Delta_{\Delta_m}^\alpha \psi(\mathbf{x}, t) + V(\mathbf{x})\psi(\mathbf{x}, t) - \kappa |\psi(\mathbf{x}, t)|^{2p}\psi(\mathbf{x}, t),
\end{equation}
where $M > 0$ is the effective mass, $V(\mathbf{x})$ is an external real-valued potential, $\kappa \in \mathbb{R}$ is the non-linear coupling constant ($\kappa > 0$ focusing, $\kappa < 0$ defocusing), and $p > 0$ is the non-linearity exponent.

\subsection{Fundamental Conservation Laws}

\begin{theorem}[Global Conservation of Mass and Energy]
Let $\psi(\mathbf{x}, t)$ be a smooth solution to \eqref{eq:nlse_main} decaying sufficiently fast at spatial infinity. Then:
\begin{enumerate}
    \item \textbf{Conservation of Total Mass (Particle Number):}
    \begin{equation}
    \mathcal{N}[\psi(t)] \coloneqq \int_{\mathbb{R}^{m-1}} |\psi(\mathbf{x}, t)|^2 \, d\mathbf{x} = \mathcal{N}[\psi(0)].
    \end{equation}
    \item \textbf{Conservation of Simplicial Energy (Hamiltonian):}
    \begin{align}
    \mathcal{E}[\psi(t)] &\coloneqq \frac{\hbar^2}{2M}\int_{\mathbb{R}^{m-1}} \sigma_{\Delta_m}^\alpha(\mathbf{k}) |\widehat{\psi}(\mathbf{k}, t)|^2 d\mathbf{k} + \int_{\mathbb{R}^{m-1}} V(\mathbf{x})|\psi(\mathbf{x}, t)|^2 d\mathbf{x} \nonumber\\
    &\quad - \frac{\kappa}{p+1}\int_{\mathbb{R}^{m-1}} |\psi(\mathbf{x}, t)|^{2p+2} d\mathbf{x} = \mathcal{E}[\psi(0)].
    \end{align}
\end{enumerate}
\end{theorem}

\begin{proof}
\leavevmode
\begin{enumerate}
    \item Multiplying \eqref{eq:nlse_main} by $\psi^*$ and taking the imaginary part:
    \begin{equation}
    \frac{\hbar}{2}\frac{\partial |\psi|^2}{\partial t} = \operatorname{Im}\left(-\frac{\hbar^2}{2M}\psi^* \Delta_{\Delta_m}^\alpha \psi\right).
    \end{equation}
    Integrating over $\mathbb{R}^{m-1}$ and applying the self-adjointness of $\Delta_{\Delta_m}^\alpha$:
    \begin{equation}
    \frac{d}{dt}\mathcal{N}[\psi(t)] = -\frac{\hbar}{M}\operatorname{Im}\langle \psi, \Delta_{\Delta_m}^\alpha \psi \rangle_{L^2} = 0,
    \end{equation}
    since $\langle \psi, \Delta_{\Delta_m}^\alpha \psi \rangle_{L^2} \in \mathbb{R}$.
    
    \item Differentiating $\mathcal{E}[\psi(t)]$ with respect to $t$:
    \begin{equation}
    \frac{d\mathcal{E}}{dt} = 2\operatorname{Re}\int_{\mathbb{R}^{m-1}} \psi_t^* \left[ -\frac{\hbar^2}{2M}\Delta_{\Delta_m}^\alpha \psi + V(\mathbf{x})\psi - \kappa |\psi|^{2p}\psi \right] d\mathbf{x}.
    \end{equation}
    Substituting the equation of motion $i\hbar \psi_t = \mathcal{H}\psi$:
    \begin{equation}
    \frac{d\mathcal{E}}{dt} = 2\operatorname{Re}\int_{\mathbb{R}^{m-1}} \left(\frac{1}{i\hbar}\mathcal{H}\psi\right)^* (\mathcal{H}\psi) \, d\mathbf{x} = 2\operatorname{Re}\left(-\frac{1}{i\hbar}\|\mathcal{H}\psi\|_{L^2}^2\right) = 0.
    \end{equation}
\end{enumerate}
\end{proof}

\subsection{Modulational Instability on Simplicial Torus}

Consider a uniform plane-wave background state with amplitude $\psi_0(\mathbf{x}, t) = \sqrt{\rho_0} e^{-i \mu_0 t / \hbar}$, where the chemical potential is $\mu_0 = -\kappa \rho_0^p$. We perturb the background as
\begin{equation}
\psi(\mathbf{x}, t) = \left(\sqrt{\rho_0} + \delta u(\mathbf{x}, t) + i \delta v(\mathbf{x}, t)\right) e^{-i \mu_0 t / \hbar}.
\end{equation}

\begin{theorem}[Simplicial Modulational Instability]
Linearizing \eqref{eq:nlse_main} with respect to perturbations $(\delta u, \delta v) \propto e^{i(\mathbf{k}\cdot\mathbf{x} - \Omega t)}$, the perturbation frequency satisfies the dispersion relation:
\begin{equation}
\hbar^2 \Omega^2(\mathbf{k}) = \frac{\hbar^2 \sigma_{\Delta_m}^\alpha(\mathbf{k})}{2M} \left[ \frac{\hbar^2 \sigma_{\Delta_m}^\alpha(\mathbf{k})}{2M} - 2\kappa p \rho_0^p \right].
\end{equation}
In the focusing regime ($\kappa > 0$), the plane wave is modulationally unstable ($\Omega^2 < 0$) for wavevectors satisfying
\begin{equation}
\sigma_{\Delta_m}^\alpha(\mathbf{k}) < \frac{4M\kappa p \rho_0^p}{\hbar^2},
\end{equation}
with maximum instability growth rate $\gamma_{\max} = \frac{\kappa p \rho_0^p}{\hbar}$ attained on the simplicial resonant manifold.
\end{theorem}

\subsection{Simplicial Solitary Waves and Anisotropic Solitons}

For stationary solitary wave states $\psi(\mathbf{x}, t) = e^{-i\mu t/\hbar} \phi(\mathbf{x})$ with $\mu < 0$:
\begin{equation}
\label{eq:soliton_profile}
-\frac{\hbar^2}{2M}\Delta_{\Delta_m}^\alpha \phi(\mathbf{x}) + (V(\mathbf{x}) - \mu)\phi(\mathbf{x}) = \kappa |\phi(\mathbf{x})|^{2p}\phi(\mathbf{x}).
\end{equation}

\begin{theorem}[Symmetry of Simplicial Solitons]
Unlike standard fractional solitons that are radially symmetric under $O(m-1)$, the ground state solutions $\phi(\mathbf{x})$ of \eqref{eq:soliton_profile} are invariant under the discrete permutation and reflection group of the simplex $S_m$. Along the principal axes $\mathbf{e}_j - \mathbf{e}_k$, the decay profile exhibits facet-aligned algebraic directional decay rates.
\end{theorem}

% ==============================================================================
\section{Anomalous Diffusion in Anisotropic Porous Media}
% ==============================================================================

\subsection{Model Formulation}

In heterogeneous porous and fractured geological media, microscopic transport is characterized by: (i) temporal retention in pore dead-ends (subdiffusion), and (ii) anisotropic spatial hopping constrained by simplicial pore geometries.

\begin{definition}[Space-Time Fractional Simplicial Diffusion Equation]
Let $u(\mathbf{x}, t)$ denote the solute concentration. The governing equation is
\begin{equation}
\label{eq:fract_diff}
\frac{\partial^\beta u(\mathbf{x}, t)}{\partial t^\beta} = \mathcal{K}_{\text{diff}} \Delta_{\Delta_m}^\alpha u(\mathbf{x}, t), \quad \mathbf{x} \in \mathbb{R}^{m-1}, \; t > 0,
\end{equation}
with initial condition $u(\mathbf{x}, 0) = u_0(\mathbf{x})$, where $\beta \in (0, 1]$ is the temporal order of the Caputo fractional derivative:
\begin{equation}
\frac{\partial^\beta u}{\partial t^\beta}(t) \coloneqq \frac{1}{\Gamma(1-\beta)}\int_0^t (t-\tau)^{-\beta} \frac{\partial u}{\partial \tau}(\tau) \, d\tau.
\end{equation}
\end{definition}

\subsection{Exact Propagator via Mittag-Leffler Functions}

\begin{theorem}[Closed-Form Solution in Frequency Domain]
The spatial Fourier transform of the exact solution to \eqref{eq:fract_diff} is given by
\begin{equation}
\widehat{u}(\mathbf{k}, t) = E_\beta\left(-\mathcal{K}_{\text{diff}}\,\sigma_{\Delta_m}^\alpha(\mathbf{k})\,t^\beta\right) \widehat{u}_0(\mathbf{k}),
\end{equation}
where $E_\beta(z) \coloneqq \sum_{n=0}^\infty \frac{z^n}{\Gamma(\beta n + 1)}$ is the Mittag-Leffler function.
\end{theorem}

\begin{proof}
Applying the Laplace transform $\mathcal{L}\{f\}(s) = \int_0^\infty e^{-st}f(t)dt$ to the Caputo derivative yields $\mathcal{L}\{\partial_t^\beta u\}(s) = s^\beta \widetilde{u}(s) - s^{\beta-1}u(0)$. Applying the spatial Fourier transform $\mathbf{x} \to \mathbf{k}$:
\begin{equation}
s^\beta \widehat{\widetilde{u}}(\mathbf{k}, s) - s^{\beta-1}\widehat{u}_0(\mathbf{k}) = -\mathcal{K}_{\text{diff}}\,\sigma_{\Delta_m}^\alpha(\mathbf{k})\,\widehat{\widetilde{u}}(\mathbf{k}, s).
\end{equation}
Solving algebraically for $\widehat{\widetilde{u}}(\mathbf{k}, s)$:
\begin{equation}
\widehat{\widetilde{u}}(\mathbf{k}, s) = \frac{s^{\beta-1}}{s^\beta + \mathcal{K}_{\text{diff}}\,\sigma_{\Delta_m}^\alpha(\mathbf{k})} \widehat{u}_0(\mathbf{k}).
\end{equation}
Applying the inverse Laplace transform using the standard identity $\mathcal{L}^{-1}\left\{\frac{s^{\beta-1}}{s^\beta + a}\right\}(t) = E_\beta(-a t^\beta)$ immediately yields the result.
\end{proof}

\subsection{Mean Squared Displacement (MSD) and Cartan Permeability Tensor}

\begin{theorem}[Anisotropic Mean Squared Displacement]
\label{thm:msd}
Assume a localized initial condition $u_0(\mathbf{x}) = \delta(\mathbf{x})$. The first and second moments of the concentration profile satisfy:
\begin{enumerate}
    \item \textbf{Mean Drift:} $\langle \mathbf{x}(t) \rangle = \mathbf{0}$.
    \item \textbf{Mean Squared Displacement Covariance Tensor:}
    \begin{equation}
    \langle \mathbf{x} \mathbf{x}^T \rangle(t) \coloneqq \int_{\mathbb{R}^{m-1}} \mathbf{x}\mathbf{x}^T u(\mathbf{x}, t)\,d\mathbf{x} = \frac{\mathcal{K}_{\text{diff}}}{m \alpha \, \Gamma(\beta + 1)} \mathbf{A}_{m-1} \cdot t^\beta.
    \end{equation}
\end{enumerate}
\end{theorem}

\begin{proof}
Using moment-generating properties of the Fourier transform:
\begin{equation}
\langle x_j \rangle(t) = \left. i \frac{\partial \widehat{u}(\mathbf{k}, t)}{\partial k_j} \right|_{\mathbf{k}=\mathbf{0}} = \left. i E_\beta'\left(-\mathcal{K}_{\text{diff}}\sigma_{\Delta_m}^\alpha t^\beta\right) \left(-\mathcal{K}_{\text{diff}} t^\beta \frac{\partial \sigma_{\Delta_m}^\alpha}{\partial k_j}\right) \right|_{\mathbf{k}=\mathbf{0}} = 0,
\end{equation}
since $\left.\nabla_\mathbf{k}\sigma_{\Delta_m}^\alpha(\mathbf{k})\right|_{\mathbf{k}=\mathbf{0}} = \mathbf{0}$ by reflection symmetry of the regular simplex.

For the second moment tensor:
\begin{equation}
\langle x_i x_j \rangle(t) = \left. -\frac{\partial^2 \widehat{u}(\mathbf{k}, t)}{\partial k_i \partial k_j} \right|_{\mathbf{k}=\mathbf{0}} = \left. -E_\beta'(0) \left(-\mathcal{K}_{\text{diff}} t^\beta \frac{\partial^2 \sigma_{\Delta_m}^\alpha}{\partial k_i \partial k_j}\right) \right|_{\mathbf{k}=\mathbf{0}}.
\end{equation}
Since $E_\beta(z) = \frac{1}{\Gamma(1)} + \frac{z}{\Gamma(\beta+1)} + \mathcal{O}(z^2)$, we have $E_\beta'(0) = \frac{1}{\Gamma(\beta+1)}$. By Theorem~\ref{thm:dispersion}, the Hessian of the symbol at $\mathbf{k}=\mathbf{0}$ is:
\begin{equation}
\left. \frac{\partial^2 \sigma_{\Delta_m}^\alpha}{\partial k_i \partial k_j} \right|_{\mathbf{k}=\mathbf{0}} = \frac{1}{m\alpha}(\mathbf{A}_{m-1})_{ij}.
\end{equation}
Substituting this yields:
\begin{equation}
\langle x_i x_j \rangle(t) = \frac{\mathcal{K}_{\text{diff}}}{m\alpha \Gamma(\beta+1)} (\mathbf{A}_{m-1})_{ij} t^\beta,
\end{equation}
completing the proof.
\end{proof}

\begin{remark}[Analytical Exact Solution vs. Numerical Intractability]
Prior to this formulation, anomalous diffusion in anisotropic media with non-Euclidean triangular/tetrahedral channel networks lacked an exact analytical Green's function, relying on approximate finite-difference or finite-element meshes. Theorem~\ref{thm:msd} resolves this challenge by providing an explicit, closed-form Fourier--Mittag-Leffler propagator directly diagonalized by the $A_{m-1}$ Cartan metric.
\end{remark}

% ==============================================================================
\section{Conclusion}
% ==============================================================================

In this paper, we have formulated and analyzed the physical dynamics governed by the Fractional Simplicial Laplacian introduced in \cite{silvafilho2026simplex}. We established the conservation laws and anisotropic solitary wave solutions of the Nonlinear Simplicial Schr\"odinger Equation, and derived the exact Mittag-Leffler propagator and Cartan permeability tensor for anomalous porous diffusion. Future research will explore split-step Fourier numerical solvers and experimental transport validation in lithographic microfluidic rock networks.

\begin{thebibliography}{99}

\bibitem{silvafilho2026simplex}
R.~M. Silva-Filho,
\emph{Analytic Continuation of Pascal's Simplex: Continuous Multinomial Integrals, Polytope Boundary Recurrences, and Simplicial Fractional Calculus},
Treatise Chapter~03, Universidade Federal de Lavras, 2026.

\bibitem{podlubny1998}
I.~Podlubny,
\emph{Fractional Differential Equations: An Introduction to Fractional Derivatives, Fractional Differential Equations, to Methods of Their Solution and Some of Their Applications},
Mathematics in Science and Engineering, Vol.~198, Academic Press, San Diego, 1999,
\doi{10.1016/S0076-5392(99)X8001-5}.

\bibitem{samko1993}
S.~G. Samko, A.~A. Kilbas, and O.~I. Marichev,
\emph{Fractional Integrals and Derivatives: Theory and Applications},
Gordon and Breach Science Publishers, Yverdon, 1993, ISBN~978-2-88124-864-1.

\bibitem{metzler2000}
R.~Metzler and J.~Klafter,
\emph{The random walk's guide to anomalous diffusion: a fractional dynamics approach},
Physics Reports, \textbf{339} (2000), no.~1, 1--77,
\doi{10.1016/S0370-1573(00)00070-3}.

\bibitem{laskin2000}
N.~Laskin,
\emph{Fractional quantum mechanics and L\'evy path integrals},
Physics Letters A, \textbf{268} (2000), no.~4--6, 298--305,
\doi{10.1016/S0375-9601(00)00201-2}.

\end{thebibliography}

\end{document}
